
A layoff selects fathers, and a trait that the father passes on and that pays off more with the son's age---Section~\ref{sec:ageprofile}
shows that skill does---would produce the same profile; son fixed effects do not remove a trait whose return grows.
The designs below narrow the reading without assigning the layoff. Together they say that the cost follows a father's layoffs during his son's childhood
when they come from his own unstable employment, not from shocks to his plant or industry, and that it is not the
income the layoff costs.

\paragraph{What kind of layoff.} The cost sits in the layoffs that look like unstable employment, not in the costly
income shocks. Splitting the teen layoff share into parts that add up to it, the within-son cost is $-0.473^{***}$ for
layoffs with a small fall in family income and $-0.096$ for those with a large fall ($p = 0.079$ for equality);
$-0.338^{***}$ for layoffs followed by fewer than 500 hours of unemployment and 0.225 for longer spells ($p = 0.064$);
and $-0.268^{**}$ for sons whose father was laid off in two or more teen years against 0.295 for one ($p = 0.145$)
(panel D). Parental income at the son's 20--26, the money after the teens that the 13--19 measure leaves out, does not
move it ($-0.232^{**}$ against $-0.236^{**}$). If the cost were the income a layoff takes, the large losses and the long
spells would carry it.


\paragraph{What kind of father.} The costly layoff marks a father whose employment stays unstable. When his son is
20--24, a unit of teen layoff share goes with more years unemployed (0.289$^{***}$) and more new layoffs (0.132$^{***}$),
at equal parental income and the father's own teen-window wage and hours, and with somewhat lower earnings
($-0.193^{*}$); his employment does not change. The small-loss layoffs, which carry the son's cost, are the ones
followed by new layoffs (0.220$^{***}$, against 0.057 after a large loss) (panel A of Table~\ref{tab:layoff_mechanism}).

\paragraph{Before and after the layoff.} The fathers were not already on a falling path. In an event study around a
father's first layoff in his son's teens, with father and calendar-year fixed effects, his earnings in the six years
before are not significantly different from the year before (jointly $p = 0.20$), though they are about 0.1 higher three
to six years before; they then fall by 0.198$^{***}$ in the layoff year and remain 0.138$^{**}$ lower six years later, the
persistent loss of the displacement literature \citep{stevens1997persistent}. His employment shows no pre-trend ($p =
0.36$) (panel A of Table~\ref{tab:layoff_referee}). The same fathers' layoffs after the son is 26 carry nothing for the
son (0.063, s.e.\ 0.128); the fathers still observed then are a subsample on which the teen layoff's cost is itself
imprecise ($-0.124$, s.e.\ 0.122), so this placebo is weak.\footnote{2,238 sons whose father is a head in at least three
waves at their 27--35. Code: \texttt{108.LayoffRefereeChecks.py}.}

\paragraph{Layoffs in bad times.} Layoffs that fall when the father's sector or his state is in decline, and so say
less about him \citep{gibbons1991layoffs}, still cost the son, about half as much. Classing each teen layoff year as bad
times when national employment of the father's sector fell or his state's unemployment rate was at least a point above
its mean (53 percent of layoff years, against 44 percent of all window years), the change between 25--29 and 45--49 is
$-0.675^{**}$ ($p = 0.035$) for bad-times layoffs and $-1.400^{***}$ for the others; within the same son $-0.166$ ($p =
0.27$) and $-0.408^{***}$; the two do not differ significantly ($p = 0.26$; panel B).\footnote{With the sector's decline
alone: $-0.664^{*}$ against $-1.143^{***}$, $p = 0.66$ for equality within son. Code: \texttt{104.FatherPathRecession.py}.}
The cost is therefore not selection alone, but the layoffs that are the father's own carry more of it.

\paragraph{Shocks the father did not choose.} A national decline of the father's industry during the son's teens raises
the father's unemployment but carries no cost to the son. With sector and birth-year fixed effects, the share of the
son's 12--18 in which national employment of the father's sector fell raises the father's unemployment share by
0.080$^{**}$ ($p = 0.014$) and his layoff share by 0.030$^{*}$ ($p = 0.086$; $F = 2.9$), and the son's earnings profile
does not bend: the within-son reduced form is 0.143$^{**}$, of the wrong sign for a cost, and the placebo windows at
5--11 and 19--25 are zero (panel A of Table~\ref{tab:layoff_exog}). The first stage on layoffs is too weak to test the
layoff cost itself: the cost the layoff path would imply, $0.030 \times -0.24$, is a tenth of the reduced form's
standard error. What the design does show is that demand-driven unemployment of the father, like a plant closure, is
not passed on.\footnote{Father's sector: his 3-digit industry in the earliest wave at the son's 10--13 (85 percent) or
his first teen wave, in 15 sectors; national employment from the BLS Current Employment Statistics, 1939--2025; 3,145
sons, 632 sector $\times$ birth-year clusters. Code: \texttt{93.LayoffIndustryShocks.py}.}

\begin{table}[!htbp]
\centering
\caption{Layoffs the father did not choose, daughters, and insurance at the father's wage (PSID)}
\label{tab:layoff_exog}
\resizebox{\ifdim\width>\textwidth\textwidth\else\width\fi}{!}{%
\resizebox*{!}{\ifdim\totalheight>0.85\textheight 0.85\textheight\else\totalheight\fi}{%
\begin{tabular}{lc}
\toprule
 & \begin{tabular}[b]{@{}c@{}}Estimate\end{tabular} \\
 & (1) \\
\midrule
\multicolumn{2}{l}{\textit{A. National employment decline of the father's sector at the son's 12--18}} \\
First stage: father's layoff share on the decline & $0.030^{*}$ \\
 & (0.017) \\
First stage: father's unemployment share on the decline & $0.080^{**}$ \\
 & (0.033) \\
First stage: father's closure share on the decline & $-0.006$ \\
 & (0.014) \\
Reduced form: son's ln earnings, change 45--49 minus 25--29 & $0.241$ \\
 & (0.203) \\
Reduced form within son: decline at 12--18 $\times$ age & $0.143^{**}$ \\
 & (0.063) \\
\quad placebo: decline at 5--11 $\times$ age & $-0.058$ \\
 & (0.066) \\
\quad placebo: decline at 19--25 $\times$ age & $0.093$ \\
 & (0.061) \\
OLS within son: layoff share $\times$ age, same sons & $-0.211^{**}$ \\
 & (0.096) \\
\addlinespace
\multicolumn{2}{l}{\textit{B. Sons and daughters, one reduced specification}} \\
Sons, ln hourly wage: layoff 45--49 minus 25--29 & $-0.538^{**}$ \\
 & (0.213) \\
Sons, ln hourly wage: within person, layoff $\times$ age & $-0.086$ \\
 & (0.096) \\
Sons, ln earnings: layoff 45--49 minus 25--29 & $-0.855^{***}$ \\
 & (0.257) \\
Sons, ln earnings: within person, layoff $\times$ age & $-0.180^{*}$ \\
 & (0.108) \\
Daughters, ln hourly wage: layoff 45--49 minus 25--29 & $0.387^{*}$ \\
 & (0.214) \\
Daughters, ln hourly wage: within person, layoff $\times$ age & $0.168^{**}$ \\
 & (0.075) \\
Daughters, ln earnings: layoff 45--49 minus 25--29 & $0.439$ \\
 & (0.285) \\
Daughters, ln earnings: within person, layoff $\times$ age & $0.200$ \\
 & (0.130) \\
\addlinespace
\multicolumn{2}{l}{\textit{C. Unemployment insurance simulated at the father's wage, within son}} \\
Layoff share $\times$ replacement rate $\times$ age & $-0.218^{*}$ \\
 & (0.129) \\
\quad + layoff share $\times$ father's ln wage $\times$ age & $-0.038$ \\
 & (0.187) \\
Placebo: father's ln wage $\times$ replacement rate $\times$ age & $0.027$ \\
 & (0.032) \\
\bottomrule
\end{tabular}}}
\par\smallskip
\begin{minipage}{\textwidth}\footnotesize \textit{Notes:} Panel A: 3,145 father-present sons; decline = share of the seven years at the son's 12--18 in which national employment (BLS CES) of the father's sector fell; sector and birth-year fixed effects; standard errors clustered by sector $\times$ birth year (632 clusters); first-stage $F$ for the layoff share 2.9. Panel B: the outcome on ln parental income and the father's teen-window wage, hours, unemployment, zero hours, volatility, closure and layoff shares; father present; clustered by 1968 family. Panel C: the benefit the state's law would have paid the father (half his weekly earnings between the state's minimum and maximum), over his weekly earnings, mean over the son's 12--18 (1976 and later), standardised; within-son design of Table~\ref{tab:layoff_selection}, with state-at-17 $\times$ age. Source: \texttt{93.LayoffIndustryShocks.py}, \texttt{92.LayoffDaughters.py}, \texttt{99.SimulatedReplacementRate.py}; table built by \texttt{100.PaperBFableTables.py}. $^{***}$ $p<0.01$, $^{**}$ $p<0.05$, $^{*}$ $p<0.1$.\end{minipage}
\end{table}

\begin{table}[!htbp]
\centering
\caption{What kind of father, which layoffs, where in the son's career, and the daughters (PSID)}
\label{tab:layoff_mechanism}
\resizebox{\ifdim\width>\textwidth\textwidth\else\width\fi}{!}{%
\resizebox*{!}{\ifdim\totalheight>0.85\textheight 0.85\textheight\else\totalheight\fi}{%
\begin{tabular}{lc}
\toprule
 & \begin{tabular}[b]{@{}c@{}}Estimate\end{tabular} \\
 & (1) \\
\midrule
\multicolumn{2}{l}{\textit{A. The father at the son's 20--24, on the teen layoff share}} \\
Father unemployed (share of years) & $0.289^{***}$ \\
 & (0.061) \\
Father laid off again (share of years) & $0.132^{***}$ \\
 & (0.035) \\
Father's ln earnings & $-0.193^{*}$ \\
 & (0.117) \\
Father employed (share of years) & $0.020$ \\
 & (0.052) \\
\addlinespace
\multicolumn{2}{l}{\textit{B. Layoffs in bad times (sector decline or high state unemployment) and other layoffs, son's ln earnings}} \\
Bad-times layoffs: change 45--49 minus 25--29 & $-0.675^{**}$ \\
 & (0.320) \\
Other layoffs: change 45--49 minus 25--29 & $-1.400^{***}$ \\
 & (0.378) \\
Bad-times layoffs: within son $\times$ age & $-0.166$ \\
 & (0.149) \\
Other layoffs: within son $\times$ age & $-0.408^{***}$ \\
 & (0.139) \\
\quad equal within son: $p$ & 0.263 \\
\addlinespace
\multicolumn{2}{l}{\textit{C. The son's career: change 45--49 minus 25--29 unless stated}} \\
ln hourly wage & $-0.603^{***}$ \\
 & (0.203) \\
\quad occupation wage score (between occupations) & $-0.171^{**}$ \\
 & (0.076) \\
\quad wage within the occupation & $-0.392^{**}$ \\
 & (0.188) \\
Same-employer annual wage growth & $0.028$ \\
 & (0.064) \\
Displaced in the year & $-0.038$ \\
 & (0.067) \\
Tenure with the employer (years), at 45--49 & $-4.214^{*}$ \\
 & (2.506) \\
Good job at 25--34 (one row per son) & $-0.041$ \\
 & (0.028) \\
Within-son ln earnings slope, early job held fixed & $-0.225^{**}$ \\
 & (0.095) \\
\bottomrule
\end{tabular}}}
\par\smallskip
\begin{minipage}{\textwidth}\footnotesize \textit{Continued on the next table; notes below its last part.}\end{minipage}
\end{table}
\begin{table}[!htbp]
\centering
\addtocounter{table}{-1}
\caption{What kind of father, which layoffs, where in the son's career, and the daughters (PSID) (continued)}
\resizebox{\ifdim\width>\textwidth\textwidth\else\width\fi}{!}{%
\resizebox*{!}{\ifdim\totalheight>0.85\textheight 0.85\textheight\else\totalheight\fi}{%
\begin{tabular}{lc}
\toprule
 & \begin{tabular}[b]{@{}c@{}}Estimate\end{tabular} \\
 & (1) \\
\midrule
\multicolumn{2}{l}{\textit{D. Daughters and sons, one reduced specification}} \\
Daughters: ln wage within person, + partnered and full year & $0.169^{**}$ \\
 & (0.075) \\
Daughters: employed, within person & $-0.073$ \\
 & (0.052) \\
Daughters: partner's ln earnings, change & $0.764^{***}$ \\
 & (0.240) \\
Daughters: earnings with zeros (Poisson), change & $0.188$ \\
 & (0.418) \\
Sons: earnings with zeros (Poisson), change & $-0.523^{**}$ \\
 & (0.243) \\
\bottomrule
\end{tabular}}}
\par\smallskip
\begin{minipage}{\textwidth}\footnotesize \textit{Notes:} Coefficients of the father's teen layoff share (per unit share; one layoff year = one seventh). Panel A: one row per father-present son, with ln parental income, the father's teen-window ln wage and ln hours and cohort dummies. Panels B--C: at equal ln parental income and every component of Table~\ref{tab:passed_on} panel A; band changes with 200 common bootstrap draws, within-son slopes clustered by 1968 family. Panel B: a bad-times layoff year is one in which national employment of the father's sector fell or his state's unemployment rate was at least a point above its 1976--2023 mean. Panel C: the son's ln wage is his occupation's CPS wage score plus his wage within it; a good job is a modal occupation $\times$ industry cell at 25--34 in the well-paid, fast-growing set. Panel D: the specification of Table~\ref{tab:layoff_exog} panel B. Source: \texttt{104.FatherPathRecession.py}, \texttt{105.DaughtersMechanism.py}, \texttt{106.LayoffCareerMechanism.py}; table built by \texttt{100.PaperBFableTables.py}. $^{***}$ $p<0.01$, $^{**}$ $p<0.05$, $^{*}$ $p<0.1$.\end{minipage}
\end{table}

\begin{table}[!htbp]
\centering
\caption{Pre-trends, a later-layoff placebo, bounds, ranks, the daughters' selection and multiple testing (PSID)}
\label{tab:layoff_referee}
\resizebox{\ifdim\width>\textwidth\textwidth\else\width\fi}{!}{%
\resizebox*{!}{\ifdim\totalheight>0.85\textheight 0.85\textheight\else\totalheight\fi}{%
\begin{tabular}{lc}
\toprule
 & \begin{tabular}[b]{@{}c@{}}Estimate\end{tabular} \\
 & (1) \\
\midrule
\multicolumn{2}{l}{\textit{A. The father around his first layoff in a son's teens (event study; reference: one year before)}} \\
Father's ln earnings, event time -6 & $0.106^{*}$ \\
 & (0.059) \\
Father's ln earnings, event time -4 & $0.085$ \\
 & (0.056) \\
Father's ln earnings, event time -3 & $0.113^{**}$ \\
 & (0.050) \\
Father's ln earnings, event time -2 & $-0.005$ \\
 & (0.058) \\
Father's ln earnings, event time 0 & $-0.198^{***}$ \\
 & (0.059) \\
Father's ln earnings, event time 1 & $-0.249^{***}$ \\
 & (0.062) \\
Father's ln earnings, event time 3 & $-0.159^{***}$ \\
 & (0.052) \\
Father's ln earnings, event time 6 & $-0.138^{**}$ \\
 & (0.060) \\
\quad pre-event coefficients jointly zero: $p$ & 0.203 \\
Father employed, event time +1 & $-0.073^{**}$ \\
 & (0.031) \\
\quad pre-event coefficients jointly zero: $p$ & 0.363 \\
\addlinespace
\multicolumn{2}{l}{\textit{B. Placebo, bounds and ranks}} \\
Layoff share at the son's 27--35 $\times$ age & $0.063$ \\
 & (0.128) \\
\quad teen layoff share $\times$ age, same sons & $-0.124$ \\
 & (0.122) \\
Largest closure cost in the 95\% interval, as a share of the layoff change & 59\% \\
Minimum detectable layoff change / within-son slope & 0.683 / 0.288 \\
Rank change (percentiles), 45--49 minus 25--29 & $-22.163^{***}$ \\
 & (7.483) \\
\addlinespace
\multicolumn{2}{l}{\textit{C. The daughters: the mother, who is observed, schooling}} \\
Mother employed in the father's layoff waves (child FE) & $0.011$ \\
 & (0.010) \\
Daughters' slope $\times$ (mother worked more after the layoff) & $0.024$ \\
 & (0.113) \\
Daughters observed with a wage at 45--49, on the layoff share & $0.058$ \\
 & (0.083) \\
Daughters' within-person wage slope, inverse-probability weighted & $0.176^{**}$ \\
 & (0.076) \\
Sons' within-person wage slope, inverse-probability weighted & $-0.134$ \\
 & (0.095) \\
Years of schooling, daughters minus sons & $0.893^{*}$ \\
 & (0.491) \\
\bottomrule
\end{tabular}}}
\par\smallskip
\begin{minipage}{\textwidth}\footnotesize \textit{Continued on the next table; notes below its last part.}\end{minipage}
\end{table}
\begin{table}[!htbp]
\centering
\addtocounter{table}{-1}
\caption{Pre-trends, a later-layoff placebo, bounds, ranks, the daughters' selection and multiple testing (PSID) (continued)}
\resizebox{\ifdim\width>\textwidth\textwidth\else\width\fi}{!}{%
\resizebox*{!}{\ifdim\totalheight>0.85\textheight 0.85\textheight\else\totalheight\fi}{%
\begin{tabular}{lc}
\toprule
 & \begin{tabular}[b]{@{}c@{}}Estimate\end{tabular} \\
 & (1) \\
\midrule
\multicolumn{2}{l}{\textit{D. Multiple testing: $p$ $\to$ adjusted $p$ (Holm / Romano--Wolf)}} \\
Layoff change (family of 6) & 0.0002 $\to$ 0.0011 / 0.0030 \\
Within-son slope (family of 6) & 0.0219 $\to$ 0.0437 / 0.0350 \\
Daughters $-$ sons, wage change (family of 4) & 0.0014 $\to$ 0.0043 / 0.0180 \\
Daughters $-$ sons, wage within person & 0.0489 $\to$ 0.0609 / 0.0550 \\
\bottomrule
\end{tabular}}}
\par\smallskip
\begin{minipage}{\textwidth}\footnotesize \textit{Notes:} Panel A: father-years at 25--64 of the fathers of father-present sons; father and calendar-year fixed effects and single-age dummies; event time binned at $-6$ and $+6$; treated = the 212 fathers with a layoff in a son's teen window. Panel B: within-son design of Table~\ref{tab:layoff_selection}; minimum detectable effects for 80\% power at 5\%; ranks are weighted percentiles within single year of age $\times$ calendar year. Panel C: the reduced specification of Table~\ref{tab:layoff_exog} panel B; inverse-probability weights from a logit of being observed with a wage in the band. Panel D: families fixed before adjusting (six tests of the layoff cost; four of the sons--daughters difference); Romano--Wolf step-down from 1,000 joint cluster-bootstrap draws of 1968 families. Clustered by 1968 family. Source: \texttt{108.LayoffRefereeChecks.py}, \texttt{109.DaughtersRefereeChecks.py}, \texttt{110.MultipleTesting.py}, \texttt{114.RomanoWolf.py}; table built by \texttt{100.PaperBFableTables.py}. $^{***}$ $p<0.01$, $^{**}$ $p<0.05$, $^{*}$ $p<0.1$.\end{minipage}
\end{table}


\subsection{Which parental income is passed on}

The same contrast shows when parental income itself is split by its source. Table~\ref{tab:shea_spell} splits log
parental income into the part predicted by the parents' spells and the rest, and regresses the son's log earnings on
both parts. The spells predict 18.2 percent of the variance of parental income, and that part passes on more strongly
than the rest: 0.528$^{***}$ against 0.360$^{***}$, a difference of 0.168$^{***}$ ($p = 0.007$). The excess comes from
the selected spells. Parental income predicted by the father's layoffs passes on with an elasticity of 0.742$^{***}$
against 0.381$^{***}$ for the rest (difference 0.361$^{**}$, $p = 0.016$), and that predicted by the father's and the
mother's unemployment by 0.265$^{***}$ ($p = 0.009$) and 0.277$^{*}$ ($p = 0.056$) more than the rest. Income predicted
by plant closures passes on with an elasticity of 0.074 ($p = 0.90$), against 0.392$^{***}$ for the rest; the difference,
$-0.317$, is imprecise (s.e.\ 0.565) because closures predict 0.4 percent of the variance of parental income.


\citet{shea2000does} instruments parental income with sources that are plausibly beyond the parents' control---union
membership, industry and job loss---and finds that the income they predict has little effect on sons' earnings: money as
such matters little. The closure row agrees with him. The layoff and unemployment rows show the opposite of a small effect,
and they are the same spells that cost the son more as he ages: what passes on strongly is not the income a layoff
costs but what comes with a layoff.\footnote{The same holds for the son's hourly wage (layoffs 0.279$^{**}$, closures
$-0.455$, s.e.\ 0.486) and his family income (layoffs 0.385$^{**}$, closures $-0.037$). Code:
\texttt{74.FrictionIGEPooled.py}, \texttt{89.SheaContrastBySpell.py} (\texttt{output/shea\_contrast\_by\_spell.md}).}

